Bottom-up multi-person pose estimators assign detected keypoints to people using a grouping mechanism trained jointly with the detector.
The grouping therefore cannot be reused once the detector changes.
This work treats grouping as a separate module and measures the cost of that separation.
PG-GAT is a graph attention network with a type-pair attention bias, skeleton-relative position encoding, and same-type repulsion heads; it embeds the detected keypoints so that $k$-means recovers the people, and a small head trained on the frozen embeddings estimates the number of people.
Trained on synthetic scenes and fine-tuned on COCO, PG-GAT reaches a pose grouping accuracy of 0.9715 on annotated COCO validation keypoints and 0.9591 on HigherHRNet detections, against 0.9847 for the detector's jointly trained associative embedding on the same detections.
Applied without retraining to YOLO11x-pose and RTMO-L, the autonomous pipeline trails the native grouping of all three detectors by 0.020 to 0.022.
At the system level, the cost is 0.048 average precision on identical keypoints, and it concentrates in scenes with more people than the synthetic training data contained.
The learned grouping heads and modularity-based clustering tested here did not improve on $k$-means, while changes to the embedding produced the gains that transfer to real data, identifying the embedding as the limiting component.
